% Generated by docs/generate_user_guide.py. Do not edit directly.
% Source: utilities/descriptions.yaml: relative_model_evidence_html

Relative model evidence compares candidate parametric models fitted to the same dataset using the same preprocessing settings and IVIVC approach. The table is shown only for selected information-criterion metrics: AIC, AIC\textsubscript{c}, and/or BIC.

In this context, empirical support means support from the observed data that a model balances fit quality and complexity well within the displayed candidate set. Lower information-criterion values indicate stronger relative support. The absolute values are not quality grades, and negative values are allowed.

Table definitions: n is the number of fitted observations; k is the number of fitted model parameters; Delta is the difference between a model's criterion value and the lowest value in the displayed candidate set; weight is the normalized relative support for that model under the criterion; and evidence ratio compares the highest weight in the set with the model's weight. Delta values near 0 indicate similar support; larger delta values indicate lower support relative to the leading candidate under that criterion.

Similar empirical support does not mean the models are scientifically equivalent. Final model selection should also consider prediction error, cross-validation behavior, residual patterns, parameter uncertainty, model simplicity, and biological or degradation-mechanism plausibility. The app reports these summaries to support user review. It does not select, recommend, or validate a model.
